\chapter{Results and Discussion}

This chapter summarizes the observed outcomes of the internship initiatives and discusses business relevance, technical implications, and identified limitations.

\section{Results of Decisioning Modernization}

The modernization and automation of decisioning artifacts provided clear operational improvements. Rule changes became pipeline-driven instead of manually compiled, artifact publication became standardized through Artifactory, versioning improved traceability and rollback readiness, and consolidation reduced complexity caused by fragmented repositories. The primary gain is not only speed but reliability under repeated releases. In regulated financial systems, consistent and auditable release behavior is as critical as low deployment latency.

\section{Results of Stand-in Server Initiative}

The stand-in server initiative established a practical path toward reducing external dependency. An internal architecture for fallback decisioning was defined, technology stack decisions were aligned with enterprise compatibility, and a phased migration strategy was established to reduce operational risk. The largest value is strategic. Even before full replacement, a reliable backup service improves continuity and strengthens operational resilience during third-party outages.

\section{Results of Fraud Centralization Pipeline}

The event-driven fraud architecture produced a structured mechanism for cross-team visibility. Incoming fraud-check events can be streamed through Kafka, consumers update a centralized repository for blocked and investigation states, and post-investigation updates close the lifecycle loop for consistent status tracking. A centralized fraud state reduces risk of contradictory decisions across banking units. This is particularly important in identity theft cases where delayed synchronization can cause severe compliance and financial impact.

\section{Measurable Outcomes Summary}

\begin{table}[h]
	\centering
	\small
	\renewcommand{\arraystretch}{1.35}
	\begin{tabular}{|p{0.17\textwidth}|p{0.17\textwidth}|p{0.21\textwidth}|p{0.21\textwidth}|p{0.17\textwidth}|}
		\hline
		\textbf{Initiative} & \textbf{KPI} & \textbf{Baseline Challenge} & \textbf{Post-Implementation Outcome} & \textbf{Business Impact} \\
		\hline
		Decisioning modernization & Rule release lead time; manual steps per release; artifact traceability rate & Manual rule compile and publish process with fragmented ownership & Push-triggered automated compile and publish path established; versioned \texttt{.adb} artifacts stored in Artifactory & Faster policy rollout and stronger audit readiness \\
		\hline
		Stand-in decision server & Internal decision coverage (\%); third-party dependency ratio; continuity readiness score & High external dependency for cloud/local third-party engines and associated recurring cost & In-house backup architecture defined (Spring Boot + React/TS) with phased adoption strategy and parity-first rollout model & Improved resilience posture and long-term cost optimization path \\
		\hline
		Fraud centralization pipeline & Fraud state propagation latency; cross-team consistency rate; stale-status incidence & No unified repository for blocked/under-investigation transactions & Kafka producer-consumer flow with central repository updates and lifecycle status synchronization implemented & Faster fraud response and reduced contradictory decision risk \\
		\hline
	\end{tabular}
	\caption{Measurable outcomes framework across internship initiatives.}
\end{table}

\textbf{Recommended KPI tracking formulas for deployment dashboards:} \textbf{Rule Release Lead Time} is calculated as \textit{Artifact Published Timestamp} minus \textit{Rule Commit Timestamp}. \textbf{Internal Decision Coverage} is computed as the ratio of requests served by the stand-in to total decision requests, expressed as a percentage. \textbf{Fraud Propagation Latency} equals \textit{Central Repository Update Timestamp} minus \textit{Fraud Event Publish Timestamp}. \textbf{Cross-Team Consistency Rate} is the percentage of decisions aligned with the latest fraud status divided by total evaluated decisions.

\section{Limitations and Risks}

Although the outcomes are positive, some limitations remain. Pipeline reliability must be maintained to avoid build bottlenecks, the stand-in service requires sustained parity verification against existing engines, Kafka-based systems require disciplined handling of ordering, retries, and duplicates, and data governance is critical for sensitive fraud records and lifecycle retention.

\section{Key Learnings from Discussion}

The discussion across all workstreams highlights a common pattern: enterprise engineering success depends on combining automation, architecture discipline, and controlled rollout. High-quality software outcomes in banking are created when technical decisions are tied directly to business risk and compliance needs.